\chapter{Teoretyczne i metodologiczne aspekty analizy opóźnień w transporcie lotniczym} 

\section{Istota i klasyfikacja opóźnień lotniczych}

Transport lotniczy stanowi jeden z absolutnych fundamentów współczesnej, zglobalizowanej gospodarki, zapewniając bezprecedensową mobilność kapitału ludzkiego oraz ułatwiając zintegrowaną, międzynarodową wymianę handlową w czasie rzeczywistym. W dobie gospodarki opartej na modelu \textit{just-in-time}, sprawność łańcuchów dostaw, operacje logistyczne korporacji transnarodowych oraz mobilność kadry zarządczej zależą bezpośrednio od przewidywalności operacji lotniczych. Sektor ten generuje znaczący wkład w globalny produkt krajowy brutto, wykazując przy tym niezwykle silny efekt mnożnikowy – każda jednostka wartości dodanej wytworzona bezpośrednio w lotnictwie generuje dodatkowe zyski w sektorze turystycznym, handlowym, usługach wysokospecjalistycznych oraz infrastrukturze lokalnej \cite{faa2022economic}. 

Jednocześnie, lotnictwo cywilne to system o niezwykle wysokim stopniu złożoności operacyjnej (tzw. system socjotechniczny), w którym koordynacja tysięcy maszyn, załóg lotniczych, służb kontroli ruchu oraz personelu obsługi naziemnej musi odbywać się w sposób ciągły, pod rygorem ekstremalnych i bezwzględnych norm bezpieczeństwa. Ta wielowymiarowa współzależność sprawia, że jest to środowisko wysoce wrażliwe na wszelkiego rodzaju zakłócenia, których ostatecznym, mierzalnym i najbardziej odczuwalnym symptomem są opóźnienia lotów. 

W literaturze przedmiotu opóźnienie definiuje się najczęściej jako różnicę pomiędzy planowanym czasem rozkładowym a rzeczywistym czasem realizacji danej operacji lotniczej. Wraz z dynamicznym rozwojem technologii i cyfryzacją lotnictwa, głębokiej ewolucji uległ sposób precyzyjnego pomiaru tego zjawiska. Współcześnie dane o punktualności nie są już rejestrowane w oparciu o subiektywne raporty załóg, lecz są generowane automatycznie poprzez standard OOOI (\textit{Out, Off, On, In}), wspierany przez pokładowe systemy teleinformatyczne (takie jak ACARS - \textit{Aircraft Communications Addressing and Reporting System}). System ten z niezwykłą precyzją rozróżnia cztery krytyczne momenty każdego rejsu:
\begin{itemize}
    \item \textbf{Out-of-gate (Czas odkołowania):} moment zwolnienia hamulców postojowych i odepchnięcia statku powietrznego od rękawa lub stanowiska postojowego (tzw. \textit{off-block time});
    \item \textbf{Off-ground (Czas startu):} fizyczna chwila oderwania kół maszyny od nawierzchni pasa startowego;
    \item \textbf{On-ground (Czas lądowania):} moment fizycznego dotknięcia pasa startowego przez podwozie samolotu na lotnisku docelowym;
    \item \textbf{In-gate (Czas dokowania):} chwila całkowitego zatrzymania maszyny i zaciągnięcia hamulców na stanowisku docelowym (\textit{in-block time}).
\end{itemize}

Dzięki powyższej kategoryzacji badacze oraz analitycy rynkowi mogą wyodrębnić dwie główne płaszczyzny pomiaru: opóźnienie wylotu (\textit{departure delay}) oraz opóźnienie przylotu (\textit{arrival delay}) \cite{doganis}. Z perspektywy konsumenta (pasażera) kluczowe jest dotarcie do miejsca docelowego zgodnie z planem, dlatego linie lotnicze w swoich materiałach marketingowych skupiają się na punktualności przylotów. Jednakże z punktu widzenia efektywności zarządzania portem lotniczym, alokacji przepustowości oraz optymalizacji sieci (ang. \textit{Air Traffic Management} - ATM), to opóźnienie wylotu stanowi miarę pierwotną. Błąd wygenerowany na etapie odlotu determinuje bowiem przebieg operacji w całej przestrzeni powietrznej i obciąża kolejne węzły systemu.

Kluczowym rozróżnieniem w teorii ekonomiki transportu jest podział zakłóceń na opóźnienia pierwotne (\textit{primary delays}) oraz opóźnienia wtórne, nazywane inaczej reakcyjnymi (\textit{reactionary delays}). Opóźnienia pierwotne wynikają z bezpośrednich przyczyn przypisanych do danej, konkretnej operacji. Może to być nagłe załamanie pogody nad lotniskiem startowym, niespodziewana awaria techniczna (usterka maszyny), incydent z udziałem pasażera (np. nagłe pogorszenie stanu zdrowia na pokładzie) czy też obostrzenia nałożone w czasie rzeczywistym przez służby kontroli ruchu lotniczego. Z kolei opóźnienia wtórne to zakłócenia „odziedziczone” z poprzednich rotacji maszyny \cite{propagation2009}. Współczesny samolot komercyjny, będący niezwykle kosztownym aktywem trwałym, musi wykazywać wysoki stopień utylizacji. Aby generować zyski, maszyna przebywa w powietrzu lub w ciągłym ruchu rotacyjnym przez kilkanaście godzin na dobę, wykonując od kilku do kilkunastu rejsów. Z tego względu opóźnienie wygenerowane pierwszym, porannym lotem nie znika. Przeciwnie – potęguje się ono w miarę upływu dnia, tworząc tzw. efekt propagacji (fali), co sprawia, że operacje wieczorne są zazwyczaj obarczone największym błędem harmonogramu.

Z punktu widzenia standardów międzynarodowych organizacji lotniczych (takich jak amerykańska Federalna Administracja Lotnictwa - FAA czy europejska organizacja Eurocontrol), nie każde, nawet kilkuminutowe odchylenie od planu jest traktowane jako usterka operacyjna. Powszechnie przyjętym, globalnym standardem branżowym jest tolerancja wynosząca 15 minut \cite{eurocontrol2024}. Przedział ten (od 0 do 15 minut) uznaje się za błąd dopuszczalny, wynikający z natury procesów stochastycznych i nieuniknionych mikrozakłóceń w obsłudze naziemnej. Dopiero przekroczenie tego 15-minutowego progu uruchamia w liniach lotniczych procedury rygorystycznego raportowania przyczyn opóźnienia, co realizowane jest za pomocą ujednoliconego, globalnego systemu kodów IATA (\textit{International Air Transport Association Delay Codes}). System ten pozwala na bezbłędną identyfikację wąskiego gardła – począwszy od kodów serii 10-19 (błędy w obsłudze pasażerskiej i bagażowej), poprzez serię 40-49 (ograniczenia techniczne i awarie), aż po kody serii 80 (restrykcje nałożone przez kontrolę lotów i pogodę).

Należy w tym miejscu zwrócić uwagę na kontrowersyjne, lecz powszechnie stosowane zjawisko sztucznego polepszania statystyk przez przewoźników, znane w literaturze jako \textit{schedule padding} (buforowanie czasu rozkładowego). Ze względu na to, że czas lotu podawany pasażerom to tzw. \textit{block time} (czas od odepchnięcia do dokowania), linie lotnicze celowo wliczają w niego ukryte rezerwy czasowe \cite{doganis}. Przykładowo, lot trwający fizycznie 60 minut może zostać wpisany w rozkład jako rejs 85-minutowy. Pozwala to linii lotniczej na „zgubienie” niewielkich opóźnień wylotu jeszcze w powietrzu. Pilot, wiedząc o opóźnieniu startu, może poprosić o bezpośrednie wektorowanie (skrót trasy) lub zwiększyć prędkość przelotową (ang. \textit{cost index}), co wiąże się z większym spalaniem paliwa, ale pozwala dotrzeć do portu docelowego zgodnie z rozkładem. Proceder ten sztucznie poprawia statystyki punktualności przylotów w oczach pasażerów oraz urzędów regulacyjnych. Z tego powodu w profesjonalnych badaniach analitycznych to opóźnienie wylotu (\textit{departure delay}) uchodzi za znacznie czystszą, obiektywną miarę efektywności infrastruktury i samych przewoźników.

Opóźnienia wykraczające poza 15-minutową tolerancję charakteryzują się nieliniowym stopniem dotkliwości operacyjnej. Z tego względu w badaniach ilościowych – w szczególności tych wykorzystujących modele uogólnione – zjawisko to dzieli się najczęściej na kategoryczne interwały \cite{greene2003}:
\begin{itemize}
    \item \textbf{Opóźnienia małe i akceptowalne (do 15 minut):} Traktowane jako naturalny element środowiska operacyjnego. Nie mają systematycznego, negatywnego wpływu na siatkę połączeń, nie zaburzają pracy załóg i są w większości niwelowane w trakcie trwania rejsu dzięki sprzyjającym prądom strumieniowym lub wspomnianym rezerwom czasowym (\textit{padding});
    \item \textbf{Opóźnienia średnie (16-45 minut):} Stanowią pierwsze, poważne zagrożenie dla stabilności systemu. Naruszają one zaplanowany czas obsługi naziemnej (\textit{turnaround time}) na kolejnym lotnisku, drastycznie zwiększają ryzyko utraty gwarantowanych przesiadek przez pasażerów w portach węzłowych oraz wymuszają na personelu naziemnym pracę pod presją czasu, co zwiększa ryzyko błędów;
    \item \textbf{Opóźnienia duże i krytyczne (powyżej 45 minut):} Wywołują głębokie perturbacje operacyjne o charakterze systemowym. Zakłócenia tego rzędu nierzadko zmuszają przewoźników do podmiany załóg ze względu na bezwzględnie rygorystyczne regulacje prawne dotyczące maksymalnego czasu pracy personelu latającego (tzw. normy FTL - \textit{Flight Time Limitations}). Mogą również skutkować koniecznością uruchamiania maszyn rezerwowych, masowym przenoszeniem pasażerów do hoteli lub całkowitym odwoływaniem rotacji (efekt lawinowy).
\end{itemize}

Zjawisko braku punktualności generuje gigantyczne, wielomiliardowe obciążenia dla światowej gospodarki, które w literaturze rozpatruje się na trzech głównych płaszczyznach kosztowych \cite{doganis}. Pierwszą z nich stanowią bezpośrednie koszty operacyjne linii lotniczych (tzw. \textit{Direct Operating Costs - DOC}). Obejmują one drastycznie zwiększone zużycie paliwa lotniczego w wyniku dłuższego kołowania (oczekiwania w kolejce do pasa) lub wymuszonego krążenia w strefach holdingowych nad lotniskiem docelowym. W skład DOC wchodzą również wydatki na nadgodziny dla obsługi naziemnej i załóg, a także potężne koszty odszkodowań konsumenckich. W samej tylko Unii Europejskiej restrykcyjne Rozporządzenie WE 261/2004 nakłada na przewoźników obowiązek wypłaty zryczałtowanych odszkodowań (od 250 do 600 EUR na pasażera) za opóźnienia przekraczające 3 godziny, co stanowi istotne zagrożenie dla płynności finansowej linii lotniczych. Do tego dochodzą tzw. koszty miękkie, czyli długoterminowa utrata zaufania i lojalności korporacyjnej (ang. \textit{goodwill}) ze strony najbardziej dochodowego, wysokomarżowego segmentu klientów biznesowych.

Drugą płaszczyzną są koszty ponoszone przez samych pasażerów i gospodarkę makro, definiowane w ekonomii głównie jako koszt utraconych możliwości (\textit{opportunity cost / Value of Time - VoT}). Każda minuta spędzona przez menedżera, inżyniera czy turystę w poczekalni lotniskowej przekłada się na mierzalny spadek produktywności gospodarki. Powoduje to niemożność odbycia kluczowych spotkań handlowych, przerwanie łańcuchów logistycznych, a na poziomie jednostkowym generuje stres i zmęczenie, co pośrednio wpływa na spadek użyteczności całego transportu lotniczego w oczach społeczeństwa.

Trzeci obszar to negatywne efekty zewnętrzne (\textit{externalities}), w tym koszty ekologiczne i systemowe infrastruktury. Wydłużony czas lotów i oczekiwanie z włączonymi silnikami na drogach kołowania generują setki tysięcy ton nadmiarowej emisji dwutlenku węgla (CO2) oraz tlenków azotu (NOx). Ponadto ciągłe zatory powodują przeciążenie kontrolerów ruchu lotniczego, co obniża margines bezpieczeństwa w zarządzaniu przestrzenią powietrzną.

\section{Uwarunkowania i determinanty punktualności lotów}

Identyfikacja i dokładna kwantyfikacja przyczyn opóźnień stanowi najważniejsze wyzwanie zarówno dla zarządców infrastruktury lotniskowej, analityków danych w liniach lotniczych, jak i instytucji państwowych. Prawidłowe modelowanie tego zjawiska wymaga zrozumienia, że punktualność jest wypadkową interakcji dziesiątek zmiennych. W literaturze naukowej oraz w oficjalnych raportach agencji lotniczych uwarunkowania te dzieli się tradycyjnie na dwie nadrzędne kategorie: czynniki egzogeniczne (zewnętrzne, w dużej mierze niezależne od decydentów biznesowych) oraz endogeniczne (wewnętrzne, wynikające wprost z przyjętych procesów operacyjnych i modeli biznesowych) \cite{iata_delay_codes}.

W powszechnej świadomości pasażerów oraz w relacjach medialnych, absolutnie głównym winowajcą opóźnień są niekorzystne warunki atmosferyczne. Rzeczywiście, parametry meteorologiczne mają bezpośredni, twardy wpływ na bezpieczeństwo operacji oraz przepustowość krzyżową lotnisk (tzw. \textit{runway throughput}). Ekstremalne zjawiska, takie jak silne wiatry boczne (\textit{crosswinds}) czy gwałtowne zmiany kierunku i siły wiatru przy samej ziemi (\textit{wind shear}), wymuszają zmianę kierunków operacji na pasach startowych, co często wymaga wstrzymania ruchu na czas rekonfiguracji podejść. Gwałtowne burze konwekcyjne w miesiącach letnich generują z kolei problemy na etapie przelotu (\textit{en-route}), zmuszając samoloty do omijania stref niebezpiecznych w promieniu kilkudziesięciu mil, co wydłuża trasę i czas lotu. W miesiącach zimowych intensywne opady śniegu lub marznący deszcz uruchamiają konieczność przeprowadzenia niezwykle czasochłonnych procedur odladzania maszyn (\textit{de-icing}) na specjalnych stanowiskach przed startem, a także wymuszają ciągłe odśnieżanie dróg kołowania. Podobnie skrajnie niska widoczność (mgła) zmusza porty lotnicze do aktywacji rygorystycznych procedur LVP (\textit{Low Visibility Procedures}), które ze względów bezpieczeństwa drastycznie zwiększają wymagane separacje czasowe i odległościowe między lądującymi samolotami. Należy jednak stanowczo zaznaczyć, że postęp technologiczny – w tym zaawansowane systemy lądowania według przyrządów (ILS Kategorii III), pozwalające na lądowania bez widoczności ziemi – potrafią w znacznym stopniu mitygować te zagrożenia \cite{doganis}. Współcześni analitycy zwracają uwagę, że to, co oficjalnie raportowane jest często jako "opóźnienie pogodowe", w rzeczywistości maskuje niewydolność operacyjną przewoźnika lub brak elastyczności systemu w warunkach zwiększonego stresu operacyjnego.

Z tego powodu najnowsza literatura transportowa coraz silniej akcentuje czynniki endogeniczne, w szczególności organizacyjne, jako prawdziwie kluczowe determinanty powstawania nieodwracalnych zatorów. Najważniejszym aspektem analitycznym jest w tym kontekście model biznesowy i specyfika polityki konkretnej linii lotniczej. To właśnie decyzje podejmowane w zarządach linii determinują elastyczność i odporność siatki połączeń na ewentualne błędy \cite{doganis}. 

W ujęciu strukturalnym wyraźnie zarysowuje się konflikt między dwoma dominującymi modelami: tradycyjnymi przewoźnikami sieciowymi (\textit{Legacy Carriers}, np. Lufthansa, Delta Airlines, LOT) a tanimi liniami lotniczymi (\textit{Low-Cost Carriers - LCC}, np. Ryanair, Wizz Air, Spirit). Przewoźnicy niskokosztowi budują swoją agresywną rentowność rynkową na maksymalizacji stopnia wykorzystania floty (\textit{aircraft utilization rate}). Aby samolot przynosił zyski przy bardzo niskich cenach biletów, musi latać nieprzerwanie, nawet do 14-15 godzin na dobę. Wymusza to na liniach LCC planowanie ekstremalnie krótkich czasów obsługi maszyny na ziemi, czyli tzw. \textit{turnaround time}. Pomiędzy dojazdem samolotu na stanowisko a jego ponownym wypchnięciem mija często zaledwie od 25 do 35 minut \cite{lcc_patterns}. Taka mordercza optymalizacja sprawia, że system zostaje całkowicie pozbawiony jakichkolwiek rezerw i buforów. Nawet najmniejsze, kilkuminutowe potknięcie – spóźniony pasażer, zacięty rękaw lotniskowy, problem z tankowaniem czy wolniejsze sprzątanie kabiny – staje się matematycznie niemożliwe do nadrobienia. W rezultacie opóźnienie zyskuje charakter chroniczny i przenosi się na każdy kolejny lot danej maszyny aż do późnej nocy. 

Z kolei potężni tradycyjni przewoźnicy sieciowi zmagają się z zupełnie innym wyzwaniem strukturalnym – koniecznością synchronizacji wielkich fal przesiadkowych. Ich model operacyjny oparty jest na scentralizowanej architekturze piasty i szprych (\textit{hub-and-spoke}). Duże lotnisko węzłowe (np. Frankfurt, Atlanta) gromadzi ruch z mniejszych lotnisk regionalnych (\textit{feeder flights}), aby połączyć pasażerów na zyskowne, międzykontynentalne trasy długodystansowe. Model ten wymusza na dyspozytorach linii lotniczych podejmowanie skrajnie trudnych decyzji. W sytuacji, gdy jeden mały samolot dowożący kilkunastu kluczowych pasażerów transferowych ulega opóźnieniu, przewoźnik celowo wstrzymuje wylot wielkiego samolotu transoceanicznego \cite{propagation2009}. Taka polityka oczekiwania ratuje wizerunek firmy u klienta i pozwala uniknąć wypłaty gigantycznych odszkodowań czy rezerwacji hoteli, ale bezlitośnie i świadomie pogarsza ogólne statystyki punktualności wyjściowej hubu, opóźniając start o kilkadziesiąt minut. Mechanizm ten generuje najgroźniejszy w lotnictwie efekt domina (\textit{delay propagation}), w którym jedno zakłócenie na skraju sieci, wskutek zależności kapitałowych i ludzkich, w ułamku doby paraliżuje operacje na zupełnie innym kontynencie.

Nie bez znaczenia dla holistycznego modelowania punktualności pozostają również uwarunkowania infrastrukturalne oraz zmienne czasowe. System lotniczy wykazuje nieprawdopodobnie silną sezonowość. Miesiące letnie generują zakłócenia wynikające z ekstremalnego nasycenia europejskiej i amerykańskiej przestrzeni powietrznej wakacyjnymi lotami czarterowymi oraz zwiększoną aktywnością burzową. Z kolei miesiące zimowe obciążają przepustowość procedurami utrzymania nawierzchni pasów startowych. Co więcej, fizyczna przepustowość lotnisk, dróg kołowania i terminali jest wielkością stałą, podczas gdy popyt wykazuje ostre wahania w cyklu dobowym. Pasażerowie preferują wyloty w określonych porach dnia, co sprawia, że harmonogramy lotów grupują się w potężne fale wylotowe i przylotowe (\textit{waves / banks / peak hours}). Poranne operacje zazwyczaj startują z "czystą kartą", a system posiada jeszcze wolne zasoby, co skutkuje najwyższą punktualnością. Jednakże w miarę upływu dnia i nawarstwiania się opisanych wyżej problemów linii LCC oraz ograniczeń sieciowych hubów, systematycznej degradacji ulega płynność całego ruchu. Pora dnia staje się zatem kluczową zmienną kontrolną, determinującą prawdopodobieństwo natrafienia na załamanie przepustowości. 

\section{Metody ilościowe w badaniach transportu lotniczego}

Współczesny globalny transport lotniczy to system generujący każdego dnia dziesiątki petabajtów danych. Powszechna cyfryzacja zarządzania ruchem lotniczym, obligatoryjna implementacja systemów ciągłego śledzenia statków powietrznych opartych na satelitarnej technologii ADS-B (\textit{Automatic Dependent Surveillance–Broadcast}) oraz funkcjonowanie otwartych, rządowych repozytoriów danych (takich jak gigantyczne bazy amerykańskiego \textit{Bureau of Transportation Statistics} - BTS) sprawiły, że branża lotnicza stała się jednym z najbardziej rozbudowanych i wymagających środowisk analitycznych typu Big Data. Ze względu na ten hipermasowy i wielowymiarowy charakter zbiorów danych, rzetelne badanie zjawiska opóźnień wykracza daleko poza proste narzędzia statystyki opisowej. Identyfikacja rzeczywistych zależności i predykcja zdarzeń wymusza na badaczach zjawisk transportowych stosowanie coraz bardziej zaawansowanego, nieliniowego aparatu matematycznego.

Historycznie w badaniach nad rynkiem transportowym absolutnie dominowała klasyczna, strukturalna ekonometria. Pierwotne próby modelowania punktualności traktowały opóźnienie wprost jako zmienną ciągłą (wyrażoną w dokładnych minutach). Do szacowania wpływu czynników (takich jak odległość trasy czy temperatura) wykorzystywano najczęściej Klasyczną Metodę Najmniejszych Kwadratów (regresja wieloraka KMNK / OLS). Narzędzie to, choć stanowi kamień węgielny klasycznej ekonomii i ekonometrii, bardzo szybko ujawniło swoje fundamentalne bariery w starciu ze specyfiką danych lotniczych \cite{greene2003}. Empiryczny rozkład opóźnień charakteryzuje się bowiem skrajną asymetrią i silną prawostronną skośnością. Dominująca większość rejsów wykonywana jest punktualnie lub z odchyleniem rzędu kilku minut (co w statystyce określa się mianem problemu nadmiaru zer – ang. \textit{zero-inflated data distribution}). Z drugiej jednak strony, rozkład ten posiada charakterystyczny „gruby ogon” (\textit{fat tail}), w którym kryją się stosunkowo rzadkie, ale ekstremalne i kluczowe poznawczo zatory rzędu kilkuset minut. Próba sztucznego dopasowania sztywnej, liniowej funkcji regresji do tak rozproszonych danych skutkuje kategorycznym łamaniem fundamentalnych założeń metody Gaussa-Markowa. Analitycy zmagają się wówczas z problemem braku normalności rozkładu reszt oraz drastyczną heteroskedastycznością wariancji składnika losowego. W efekcie estymatory OLS tracą cechę efektywności, a wnioskowanie statystyczne oparte na tradycyjnych testach (np. teście t-Studenta) staje się niewiarygodne i potencjalnie obciążone.

Świadomość tych ograniczeń doprowadziła do rewolucji metodycznej. Na przełomie wieków badacze zjawisk transportowych zaczęli masowo odchodzić od prób predykcji dokładnej, pojedynczej minuty spóźnienia samolotu. Zamiast tego zwrócili się ku klasie modeli zmiennych jakościowych (dyskretnych), znacznie lepiej odpowiadających biznesowym potrzebom linii lotniczych. Podejście to opiera się na transformacji ciągłego czasu trwania zdarzenia w logicznie uporządkowane szufladki kategorialne (np. opóźnienie małe, średnie, duże) i wyliczaniu prawdopodobieństwa trafienia konkretnego lotu do jednej z nich. Największą popularność na tym polu zyskały probabilistyczne modele logitowe i probitowe. 

Szczególnym i niezwykle potężnym w zastosowaniach lotniczych wariantem tego podejścia stał się Uporządkowany Model Logitowy (\textit{Ordered Logit} lub model proporcjonalnych szans - \textit{Proportional Odds Model}) \cite{greene2003}. Jego genialność polega na zdolności do matematycznego odzwierciedlenia naturalnej, rosnącej hierarchii zjawiska. Mechanika modelu uwzględnia fakt, że przejście z kategorii opóźnienia małego do średniego to postępujący ciąg tego samego zjawiska. Największą, pragmatyczną zaletą modeli z rodziny logitowej jest jednak ich transparentność i interpretowalność wyników. Wynikiem estymacji parametrów są łatwo mierzalne Ilorazy Szans (\textit{Odds Ratios - OR}). Pozwalają one badaczom i menedżerom wprost sformułować wnioski biznesowe w rodzaju: „wykonanie rejsu późnym wieczorem zwiększa szansę na wpadnięcie do gorszej kategorii opóźnienia o 45\% w stosunku do lotu porannego”. 

Ostatnia dekada przyniosła jednak w naukach o danych kolejny potężny paradygmat, który coraz silniej przenika do ekonomiki transportu – zastosowanie algorytmów uczenia maszynowego (\textit{Machine Learning - ML}). Metody bazujące na sztucznej inteligencji, takie jak Lasy Losowe (\textit{Random Forest - RF}), Maszyny Wektorów Nośnych (\textit{Support Vector Machines - SVM}), zaawansowany Gradient Boosting (np. algorytmy XGBoost) czy też Głębokich Sieciach Neuronowych (\textit{Deep Learning / ANN}), stanowią całkowicie odmienną, inżynieryjną filozofię badawczą \cite{albassam2025}. W przeciwieństwie do klasycznej ekonometrii, nowoczesne algorytmy ML wywodzą się z rodziny metod nieparametrycznych. Oznacza to, że nie wymuszają na analityku wcześniejszego określenia kształtu funkcji, nie boją się kolinearności zmiennych i nie narzucają na zbiór uczący żadnych rygorystycznych założeń dotyczących jego rozkładu czy homoskedastyczności. 

Zastosowanie algorytmów opartych na agregacji drzew decyzyjnych (\textit{ensemble learning}), takich jak Lasy Losowe, okazało się ogromnym przełomem w badaniu systemów nieliniowych. Lasy Losowe budują setki niezależnych ścieżek decyzyjnych dla różnych podzbiorów danych, a ostateczny wynik ustalają na drodze tzw. głosowania większościowego. Wykazują one wybitną, niedoścignioną dla klasycznej ekonometrii zdolność do samoczynnego wyłapywania wielowymiarowych interakcji (np. faktu, że konkretna, lokalna linia lotnicza radzi sobie świetnie w słoneczny poranek, ale całkowicie paraliżuje lotnisko podczas deszczowego popołudnia na określonym kierunku operacyjnym). Ich główną przewagą analityczną jest drastyczna poprawa zdolności klasyfikacji rzadkich anomalii systemowych (krytycznych zakłóceń >45 minut), z którymi klasyczne logity często mają problem z powodu matematycznej tendencji do uśredniania prawdopodobieństw dla klas dominujących (lotów punktualnych). Rozwój metod ML wymusił również zmianę w sposobie ewaluacji jakości modeli. Podczas gdy w regresji liniowej złotym standardem był wskaźnik R-kwadrat, w lotniczej analizie klasyfikacyjnej nadrzędną rolę zaczęły pełnić wskaźniki odporne na asymetrię danych (tzw. zbiory niezbalansowane), takie jak Precyzja, Czułość (\textit{Recall}) czy wreszcie zintegrowany wskaźnik \textit{F1-Score}.

Wybór ostatecznego narzędzia matematycznego dla celów pracy dyplomowej czy też wdrożenia rynkowego stawia analityka przed jednym z najsłynniejszych dylematów we współczesnej analityce danych: \textit{Inference vs. Prediction} (Wnioskowanie czy Prognozowanie). Badacz musi określić priorytet. Jeśli celem jest stworzenie doskonałego operacyjnie oprogramowania dla linii lotniczej, potrafiącego z najwyższą trafnością (Accuracy/F1-Score) odseparować i przewidzieć najcięższe zakłócenia, naturalnym wyborem będą czarne skrzynki uczenia maszynowego (Random Forest). Jeśli jednak nadrzędnym celem instytucji regulacyjnej (bądź autora pracy akademickiej) jest głębokie zrozumienie samej mechaniki rynkowej i ocena istotności statystycznej danego zjawiska, wybór powinien paść na rygorystyczny model ekonometryczny (Ordered Logit) \cite{albassam2025}. Modele parametryczne pozwalają otworzyć system i zajrzeć w strukturę kosztów oraz szans, podczas gdy AI dostarcza jedynie końcowy wynik. Z tego właśnie powodu w najbardziej zaawansowanych, współczesnych pracach badawczych dotyczących inżynierii transportowej kategorycznie zaleca się jednoczesne stosowanie i zderzanie ze sobą obu tych skrajnych podejść metodologicznych, co gwarantuje pełną i rzetelną triangulację wyników naukowych.